\documentclass{Thesis}
\usepackage[utf8]{inputenc}
\usepackage[T1]{fontenc}
\usepackage{geometry}
\usepackage{graphicx}
\usepackage{float}
\usepackage{hyperref}
\usepackage[english,indonesian]{babel}

\title{Optimasi Pemrosesan Template LaTeX Menggunakan Pendekatan Berbasis AI}
\author{Academic Support Agent}
\date{\today}

\begin{document}
\maketitle

\begin{abstract}

\end{abstract}

\chapter{PENDAHULUAN}


\chapter{TINJAUAN PUSTAKA}


% Original content (commented out):
% \documentclass[12pt,oneside]{Thesis}
\usepackage[spanish,mexico]{babel}
\usepackage[sectionbib, square, comma, numbers, sort&compress]{natbib} 
\hypersetup{urlcolor=black, colorlinks=true} 
\usepackage{amsmath} 
\usepackage{mathrsfs}
\usepackage{indentfirst}
\usepackage{enumitem}
\usepackage{url}
\usepackage{float}
\usepackage[colorinlistoftodos, prependcaption, textsize=tiny]{todonotes}
\setlength{\parindent}{0.5cm}
\newcommand{\tr}{\mathrm{tr}}
\newcommand{\diag}{\mathrm{diag}}
\renewcommand{\qedsymbol}{\rule{0.5em}{0.5em}}
\usepackage[spanish]{cleveref}
\usepackage[spanish]{babelbib}
\selectbiblanguage{spanish}
\usepackage{siunitx}
\usepackage{booktabs}
\usepackage{bbm}
\usepackage[many]{tcolorbox}
\usepackage[toc]{appendix}
\newtheorem{observacion}[teorema]{Observaci\'on}
\newtheorem{hipotesis}[teorema]{Hip\'otesis}
\usepackage{xpatch}
\xpatchcmd{\proof}{\itshape}{\normalfont\proofnamefont}{}{}

%%%%%%%%$%%%%%%%%
%Este pequeño bloque permite que el número de la ecuación permanezca del tamaño del resto del texto
%si se decide hacer small el tamaño de la ecuación.
%%%%%%%%%%%%%%%%%
\makeatletter
\def\maketag@@@#1{\hbox{\m@th\normalfont\normalsize#1}}
\makeatother
%%%%%%%%%%%%%%%%


\newcommand{\proofnamefont}{\bfseries}
\begin{document}
	\frontmatter 
	\setstretch{1.5} 
	\fancyhead{} 
	\rhead{\thepage} 
	\lhead{} 
	\pagestyle{fancy} 
	\newcommand{\HRule}{\rule{\linewidth}{0.5mm}} 
	\hypersetup{pdftitle={\titulo}}
	\hypersetup{pdfauthor=\autor}
%---------------------------------------------------------------------------
	\input{Portada}
	\input{Dedicatoria}
	\input{Agradecimientos}
	%\listoftodos
	\input{Contenido}
	\mainmatter 
	\pagestyle{fancy} 
  \input{Introduccion/Introduccion}
	\input{Capitulos/Capitulo1/Capitulo1}
	\input{Capitulos/Capitulo2/Capitulo2}
	\input{Capitulos/Capitulo3/Capitulo3}
	\begin{appendices}
    	\input{Apendice}
	 \end{appendices}
%---------------------------------------------------------------------------
%	BIBLIOGRAPHY
%---------------------------------------------------------------------------
	\label{Bibliography}
	\lhead{\emph{Bibliograf\'ia}}
	\bibliographystyle{plainnat} 
	\bibliography{library} 
\end{document}

\end{document}